\section{Điều hai bài chưa giải quyết}
\label{sec:ch12-dieu-chua-giai-quyet}

Mục trước khép lại ở chỗ một phép đối chiếu đã được chạy và các chỉ số không
trụ được. Cộng với mục~\ref{sec:ch12-con-so-va-cai-no-dem}, hai bài báo của
chương này đưa được ba điều, và cả ba đều là quan sát về cách một lĩnh vực
làm việc chứ không phải định lý.

Thứ nhất, khối thí nghiệm có người mà đánh giá theo chức năng vay mượn hầu như
không tồn tại: 128 bài trên 18\,254, và trong chính những bài đã tự tuyên bố
rằng phương pháp của mình giải thích được cho người thì chỉ hơn một nửa đem
lời tuyên bố ra thử. Thứ hai, con số 128 ấy là chặn trên, vì phép chấm chỉ hỏi
thí nghiệm có tồn tại hay không chứ không hỏi nó tốt hay tồi, và bài báo tự
liệt kê ba chỗ mà một thí nghiệm có thể rỗng. Thứ ba, ở lần hiếm hoi mà phép
đối chiếu được chạy đầy đủ, trên bảy chỉ số của lời giải thích phản thực, các
chỉ số không dự báo được đánh giá của người và còn đổi cả dấu giữa ba bộ dữ
liệu.

Cái hai bài không đưa được thì gom lại thành ba điều nữa.

Thứ nhất, không bài nào chạm tới các họ chỉ số của
chương~\ref{ch:do-do-trung-thuc}. Bài thứ nhất không nhắc tới độ trung thực,
tới attribution hay tới bất kỳ phương pháp nào của Phần~II; bài thứ hai chấm
các chỉ số của một họ lời giải thích khác. Cho nên chương này bổ sung một chỗ
trống nữa vào danh sách chỗ trống mà Phần~IV đang lập, chứ không lấp chỗ trống
nào.

Thứ hai, cả hai bài đều dừng ở phía tính hợp lý. Cuộc kiểm đếm đếm việc kiểm
chứng lời tuyên bố về khả năng giải thích cho người; nghiên cứu người dùng hỏi
người tham gia rằng lời giải thích có hiểu được, có hợp lý, có đủ chi tiết
không. Không câu hỏi nào trong đó là câu hỏi của định
nghĩa~\ref{def:ch01-do-trung-thuc}, và không thể là, vì người trả lời không
nhìn thấy mô hình. Mũi tên chấm vòng lên phía mô hình trong
hình~\ref{fig:ch12-hai-nhanh-danh-gia} vẫn nguyên là nét chấm sau cả hai bài
báo.

Thứ ba, và đây là chỗ mà nghiên cứu người dùng tự nhận, nó cho thấy các chỉ số
hiện có không làm được đại lượng thay thế nhưng không cho biết một chỉ số làm
được thì trông ra sao. Hướng nghiên cứu mà các tác giả nêu ở cuối bài là xây
những chỉ số đặt nền trên lý thuyết về con người rồi kiểm chứng chúng với
\tn{cảm nhận của người dùng}{user perception}~\autocite{p25userperception}. Đó là một chương trình
nghiên cứu, không phải một kết quả, và mục~\ref{sec:ch12-ket-qua-va-ranh-gioi}
vừa cho thấy chương trình ấy phải giải một bài toán mà bài báo chưa đụng tới:
cùng bảy chỉ số ấy đi ngược chiều nhau trên hai bộ dữ liệu, nên một chỉ số đặt
nền trên cảm nhận của người sẽ phải trả lời cảm nhận của người nào, ở lĩnh vực
nào.

Ngoài sáu điều trên còn một câu hỏi mà cả hai bài đều không đặt ra, và chương
này chỉ nêu chứ không trả lời được. Cả hai bài đọc tình trạng hiện nay như một
việc chưa ai làm: chưa ai chạy đủ thí nghiệm có người, chưa ai xây được chỉ số
đặt nền trên cảm nhận. Cách đọc ấy giả định rằng nếu lĩnh vực chịu bỏ công thì
việc sẽ xong. Nhưng chuỗi các chương từ
chương~\ref{ch:do-do-trung-thuc} tới đây cho thấy mỗi lần một phép kiểm chứng
được đem ra chạy thì nó lại đòi một phép kiểm chứng khác đứng dưới nó: chỉ số
đòi ground truth, ground truth dựng sẵn chỉ dựng được cho những nhiệm vụ có
cấu trúc riêng, cảm nhận của người thì đổi theo lĩnh vực. Câu hỏi là chuỗi ấy
có đáy hay không.

Chương~\ref{ch:gioi-han-ly-thuyet} nhận đúng câu hỏi ấy, và nhận nó từ một
hướng khác hẳn những chương vừa qua. Những chương ấy đi tìm bằng chứng thực
nghiệm về việc các dụng cụ đo hiện có làm được gì. Chương~\ref{ch:gioi-han-ly-thuyet}
thì hỏi một câu không cần thí nghiệm nào: liệu có một giới hạn lý thuyết cho chính khả năng giải
thích, chứng minh được bằng lý thuyết thông tin thuật toán chứ không bằng phép
đo. Nếu có một giới hạn như vậy thì phần nào của tình
trạng ở chương này là do lĩnh vực chưa làm, và phần nào là do không làm được.

\index{đánh giá!hai nhánh|)}%
